\section*{Capítulo \ref{ch:b3:quantum-angular-momentum} --- Momento angular quântico}

\begin{solution}{exo:b3:quantum-angular-momentum:1}
(a) $[\hat y\hat p_z - \hat z\hat p_y,\ \hat z\hat p_x - \hat x\hat
p_z]$: só sobrevivem os termos que compartilham um par $\hat z, \hat p_z$,
o que dá $\iu\hbar(\hat x\hat p_y - \hat y\hat p_x)$. (b) $[\hat
L^2, \hat L_z] = \sum_i[\hat L_i^2, \hat L_z] = \hat L_x[\hat L_x,
\hat L_z] + [\hat L_x, \hat L_z]\hat L_x + (x \to y)$: os quatro
termos se cancelam aos pares. (c) $[\hat L_z, \hat x] = \iu\hbar\hat y$:
$\hat L_z$ gera rotações em torno de $z$ --- tanto de operadores
quanto de estados. (d) Que duas componentes sejam simultaneamente bem
definidas obriga (pelo \omterm{def:b3:quantum-formalism:commutator}{comutador}) a terceira a ser bem definida, com
valor nulo para as três: só $l = 0$ consegue isso.
\end{solution}

\begin{solution}{exo:b3:quantum-angular-momentum:2}
(a) $\hat L_z = \hbar\operatorname{diag}(1, 0, -1)$. (b) $\hat L_+ =
\hbar\sqrt2\,\begin{pmatrix}0&1&0\\0&0&1\\0&0&0
\end{pmatrix}$, e $\hat L_-$ é sua transposta. (c) $\hat L_x =
\dfrac{\hbar}{\sqrt2}\begin{pmatrix}0&1&0\\1&0&1\\0&1&0
\end{pmatrix}$: polinômio característico $\lambda(\lambda^2 -
\hbar^2)$. (d) O autovetor de $L_x = +\hbar$ é $\tfrac12(1,
\sqrt2, 1)$: probabilidades $\tfrac14, \tfrac12, \tfrac14$ para $m =
+1, 0, -1$.
\end{solution}

\begin{solution}{exo:b3:quantum-angular-momentum:3}
(a) $I = \mu r_0^2 = \qty{1.46e-46}{kg.m^2}$; $B = \hbar^2/2I =
\qty{3.8e-23}{J} = \qty{0.24}{meV}$. (b) $2B/h = \qty{115}{GHz}$,
$\lambda = \qty{2.6}{mm}$. (c) \qty{231}{GHz} e \qty{346}{GHz}.
(d) O espaçamento dá $B$ e, com ele, $I = \mu r_0^2$: o comprimento de
ligação de uma molécula a anos-luz de distância, com aproximadamente
metade da precisão relativa do espaçamento.
\end{solution}

\begin{solution}{exo:b3:quantum-angular-momentum:4}
(a) $m\hbar$ com $m = -2 \dots 2$. (b) $\cos\theta = 2/\sqrt6$:
$\theta = 35.3^\circ$. (c) O alinhamento perfeito tornaria $L_x = L_y
= 0$ bem definidos junto com $L_z$ --- o que os \omterm{def:b3:quantum-formalism:commutator}{comutadores} proíbem,
exceto no caso trivial $l = 0$. (d) $l/\sqrt{l(l+1)} \to 1$: para $l$
grande, o cone se fecha sobre o eixo e a seta clássica retorna.
\end{solution}

\begin{solution}{exo:b3:quantum-angular-momentum:5}
(a) Sem dependência em $\varphi$, o operador dá
$-\hbar^2(\sin\theta)^{-1}\partial_\theta(\sin\theta\,(-\sin\theta))
= 2\hbar^2\cos\theta$. (b) O mesmo cálculo com o fator
$\eu^{\pm\iu\varphi}$: \omterm{def:b3:quantum-formalism:observable}{autovalor} $2\hbar^2$; $\hat L_z$ dá $\pm\hbar$.
(c) $\cos\theta$ e $\sin\theta\,\eu^{\pm\iu
\varphi}$ mudam de sinal sob $\vect r \to -\vect r$ ($\theta \to \pi
- \theta$, $\varphi \to \varphi + \pi$): paridade $-1 = (-1)^1$. (d)
Eles são autovetores do \omterm{def:b3:quantum-formalism:observable}{operador hermitiano} $\hat L_z$ com \omterm{def:b3:quantum-formalism:observable}{autovalores}
diferentes.
\end{solution}

\begin{solution}{exo:b3:quantum-angular-momentum:6}
(a) $2J + 1$ estados dividem o nível (degenerescência); o fator de
Boltzmann tributa a energia dele. (b) Maximize o produto: $\dd/\dd J = 0$
dá $2 = (2J+1)^2B/k_{\text{B}}T$, o $J_{\max}$ enunciado. (c) A
\qty{10}{K}: $J_{\max} \approx 0.8$, de modo que $J = 1$ domina e
brilham as linhas de \qty{115}{GHz} e \qty{230}{GHz}; a \qty{300}{K}:
$J_{\max} \approx 7$. (d) $T \approx 2B(J_{\max} + \tfrac12)^2/
k_{\text{B}} \approx \qty{310}{K}$ --- um termômetro rotacional.
\end{solution}

\begin{solution}{exo:b3:quantum-angular-momentum:7}
(a) O nível superior se desdobra em cinco e o inferior em três, todos
com o mesmo espaçamento $\mu_{\text{B}}B$. (b) Com espaçamentos iguais,
$h\nu = h\nu_0 + (\Delta m)\mu_{\text{B}}B$ e $\Delta m \in \{0,
\pm1\}$: apenas três posições. (c) $\mu_{\text{B}}/h = \qty{14}{GHz/T}$, de modo que em $B = \qty{2}{T}$
o desdobramento vale \qty{28}{GHz}, isto é, $\Delta\lambda =
\lambda^2\Delta\nu/c \approx \qty{23}{pm}$ --- resolvível desde as
redes de difração de Zeeman, em 1896.
(d) O spin do próprio elétron, com seu fator anômalo $g \approx
2$: os padrões ``anômalos'' foram as impressões digitais do spin antes
que o spin tivesse nome (\cref{ch:b3:spin-two-level}).
\end{solution}

\begin{solution}{exo:b3:quantum-angular-momentum:8}
Nas unidades enunciadas, $U_{\text{eff}} = -2/x + l(l+1)/x^2$. (a) Para
$l = 1$: mínimo em $x = 2$ ($r = 2a_0$), de profundidade
$-E_{\text{I}}/2$. (b) $U_{\text{eff}} > 0$ para $x < 1$: dentro de um
raio de Bohr, a muralha vence. (c) $\psi \sim r^l$ perto da origem: só
os estados $l = 0$ têm densidade não nula em $r = 0$. (d) Capturar um
elétron exige densidade eletrônica \emph{no} núcleo: os elétrons $s$,
os únicos não barrados pela muralha centrífuga, são os engolidos.
\end{solution}

\begin{solution}{exo:b3:quantum-angular-momentum:9}
(a) $\hat L_x = (\hat L_+ + \hat L_-)/2$ muda $m$: os elementos
diagonais se anulam. (b) Por simetria, os dois quadrados transversais
são iguais, e sua soma vale $\langle\hat L^2 - \hat L_z^2\rangle$. (c)
Em $m = l$: $\Delta L_x\Delta L_y = l\hbar^2/2 = \tfrac\hbar2
\,l\hbar$: igualdade --- o estado esticado é tão alinhado quanto a
álgebra permite. (d) O cone: $L_z$ bem definido, médias transversais
nulas, dispersões transversais iguais --- um vetor de comprimento e
projeção definidos, democraticamente borrado em azimute.
\end{solution}

\begin{solution}{exo:b3:quantum-angular-momentum:10}
(a) Contabilidade de energia com $E_J = BJ(J+1)$: $J \to J + 1$
acrescenta $2B(J+1)$; $J \to J - 1$ subtrai $2BJ$. (b) $\Delta J = 0$ é
proibido, e os dois ramos começam em $\pm 2B$: nada cai na frequência
vibracional pura. (c) Dois pentes de espaçamento $2B$ ladeando uma
lacuna, com intensidades que crescem até $J_{\max}$ e depois caem. (d)
A largura da banda é a extensão povoada dos pentes, que cresce como
$\sqrt{T}$: essas asas termicamente povoadas são precisamente onde uma
banda de efeito estufa saturada continua absorvendo.
\end{solution}

\begin{solution}{exo:b3:quantum-angular-momentum:11}
(a) $\|\hat J_\pm\ket{j,m}\|^2 = \bra{j,m}\hat J_\mp\hat J_\pm
\ket{j,m} = \hbar^2[j(j+1) - m(m\pm1)]$. (b) Subir além de $m = j$
daria $j(j+1) - j(j+1) = 0$ e, um passo adiante, normas negativas: a
cadeia precisa conter o estado do topo, que é aniquilado. (c) O topo e
o fundo diferem por um número inteiro de passos unitários: $2j \in
\mathbb N$, ou seja, $j = 0, \tfrac12, 1, \tfrac32, \dots$ (d) A
álgebra nunca usa funções de onda; só a exigência de que
$\eu^{\iu m\varphi}$ seja unívoca no círculo obriga $m$ a ser inteiro
--- condição que prende o movimento orbital e deixa os momentos
angulares intrínsecos (de spin) livres para serem semi-inteiros.
\end{solution}

\begin{solution}{exo:b3:quantum-angular-momentum:12}
(a) $n\hbar$ por unidade de volume. (b) $\omega = n\hbar V/(\tfrac12 Ma^2)
= 2n\hbar/\rho a^2 = 2 \times \num{8.5e28} \times \num{1.055e-34}/
(7900 \times \num{e-4}) \approx \qty{2.3e-5}{rad/s}$ --- invisível de
uma só vez. (c) Inverter em compasso com a ressonância do fio acumula
$\sim Q$ empurrões: $\omega \sim \qty{2e-3}{rad/s}$, oscilações de
miliradianos com período de segundos --- visíveis com um espelho e um
feixe de luz, como Einstein e de Haas viram. (d) A razão giromagnética
medida foi $e/m_{\text{e}}$, o dobro do valor orbital
$e/2m_{\text{e}}$: o magnetismo do ferro não vem de correntes orbitais,
mas do \emph{spin} do elétron, cujo $g \approx 2$ os próximos capítulos
explicam.
\end{solution}

\begin{solution}{pb:b3:quantum-angular-momentum:1}
\textbf{1.} $E_J = BJ(J+1)$, com degenerescência $2J + 1$.
\textbf{2.} $I = \mu r_0^2 = \qty{1.46e-46}{kg.m^2}$: $B =
\hbar^2/2I = \qty{3.8e-23}{J}$, isto é, $B/h = \qty{57.7}{GHz}$ --- a
constante medida.
\textbf{3.} $115$, $231$ e \qty{346}{GHz}.
\textbf{4.} $E_{J+1} - E_J = 2B(J+1)$ cresce linearmente: linhas
consecutivas diferem pela constante $2B$. O espaçamento dá $I$ e, com
ele, o comprimento de ligação --- química estrutural por rádio.
\textbf{5.} $\lambda = c/\nu = \qty{2.6}{mm}$. O vapor de água
atmosférico devora as ondas milimétricas; daí o Atacama, o Mauna Kea e
o Polo Sul --- altos, frios e secos.
\textbf{6.} O fóton é uma partícula de spin $1$ que carrega $\hbar$: a
conservação do momento angular total obriga o $J$ da molécula a mudar
de uma unidade.
\textbf{7.} Sua distribuição de carga é simétrica em todo ângulo de
rotação: nenhum dipolo oscilante, nenhuma linha dipolar --- camuflagem
perfeita.
\textbf{8.} $B_{\text{H}_2} = \hbar^2/2\mu r_0^2 \approx
\qty{7.6}{meV}$; as transições quadrupolares fracas da molécula
simétrica exigem $\Delta J = 2$, ao custo de $6B \approx
\qty{46}{meV}$ --- cinquenta vezes o $k_{\text{B}}T$ a \qty{10}{K}: o
hidrogênio frio não consegue irradiar de modo algum.
\textbf{9.} $2B = \qty{0.48}{meV} < k_{\text{B}}T = \qty{0.86}{meV}$:
as colisões mantêm povoados os primeiros degraus --- a escada funciona
exatamente onde as nuvens vivem.
\textbf{10.} $n_1/n_0 = 3\,\eu^{-0.478/0.862} = 1.7$: o nível emissor
está bem abastecido.
\textbf{11.} $J_{\max} \approx \sqrt{0.86/0.48} - 0.5 \approx 0.8$: a
população tem pico em $J = 1$.
\textbf{12.} O CO entrega os fótons; convertê-los em massa total de gás
repousa sobre uma abundância CO/H$_2$ suposta --- mensageiro luminoso,
resgate calibrado.
\textbf{13.} $\Delta\nu = \qty{115}{MHz}$: $v = c\,\Delta\nu/\nu
\approx \qty{300}{km/s}$, afastando-se (frequência reduzida) ---
velocidades orbitais galácticas.
\textbf{14.} Os $\pm\qty{2}{km/s}$ cobrem \qty{1.5}{MHz}; o movimento
térmico a \qty{10}{K} dá apenas $\sim\qty{40}{kHz}$: a largura é
turbulência, não temperatura --- as larguras de linha mapeiam o tempo
interno de uma nuvem.
\textbf{15.} Em torno de uma massa central, $v \propto 1/\sqrt r$
deveria cair (Kepler); a planura medida significa que a massa
englobada continua crescendo com o raio --- matéria invisível
envolvendo o disco luminoso: matéria escura.
\textbf{16.} A razão de populações, isto é, a temperatura de excitação
do gás.
\textbf{17.} Em \qty{57.6}{GHz}: a expansão cósmica dobrou todo
comprimento de onda em voo --- o desvio para o vermelho cosmológico ao
qual o último capítulo deste livro volta.
\textbf{18.} Fluxo de fótons $\to$ luminosidade em CO (é preciso a
distância) $\to$ contagem de CO (modelo de excitação) $\to$ massa de
H$_2$ (o ``fator X'' de CO para H$_2$: o elo mais fraco, calibrado
localmente e exportado com cautela).
\textbf{19.} Um núcleo a \qty{10}{K} não emite nada que um telescópio
óptico possa ver; suas linhas rotacionais são, ao mesmo tempo, seu
único termômetro, velocímetro e balança.
\textbf{20.} Um termômetro é sensível onde o espaçamento entre níveis é $\sim
k_{\text{B}}T$: um espaçamento de \qty{5.5}{K} responde
fortemente entre $5$ e $50$ K, ao passo que um nível óptico de
\qty{2}{eV} não seria povoado por nada.
\textbf{21.} A linha satura (fica opticamente espessa) e, pior, o CO
se congela sobre os grãos de poeira no gás mais denso e mais frio: os
astrônomos passam a isotopólogos mais raros ($^{13}$CO, C$^{18}$O) e a
outras moléculas.
\textbf{22.} Cada espécie precisa da sua própria densidade e
temperatura para brilhar: juntas, elas tomografam a nuvem --- os
envoltórios externos em CO, os núcleos densos em NH$_3$ e HCN.
\textbf{23.} Uma linha é vista em contraste com o fundo cósmico; à
medida que a temperatura do gás se aproxima de \qty{2.7}{K}, sua
excitação se equilibra com o fundo e o contraste desaparece: nenhuma
nuvem mais fria pode ser lida assim.
\textbf{24.} $\lambda/20 = \qty{130}{\micro m}$ contra
$\qty{25}{nm}$ do trabalho óptico: cinco mil vezes mais tolerante ---
e é por isso que antenas milimétricas de doze metros ficam ao relento
no deserto, enquanto os espelhos ópticos vivem em cúpulas.
\textbf{25.} Uma escada $BJ(J+1)$ com $B/h = \qty{57.6}{GHz}$, viva a
dez kelvins e lida em \qty{2.6}{mm}, mede a massa, a temperatura, a
turbulência e a rotação do universo frio --- e suas curvas de rotação
planas são uma exposição permanente a favor da matéria escura.
\end{solution}
